% ---------------------------------------------------------------
% preambulo.tex — pacotes, estilos e macros
%
% A classe llncs (Springer LNCS) define as margens, os tipos de letra e
% os títulos. Não se alteram aqui: mexer em margens ou tamanhos de letra
% faz o artigo deixar de cumprir o formato (e o limite de páginas).
% ---------------------------------------------------------------

\usepackage[T1]{fontenc}
\usepackage[utf8]{inputenc}
% Português. Com o ficheiro de língua clássico (portuges.ldf) usa-se esse;
% numa instalação mínima sem ele, o babel moderno carrega a língua pelo
% ficheiro .ini. O resultado é o mesmo.
% O babel clássico torna " um atalho que «come» o carácter seguinte;
% desativa-se, para que aspas escritas diretamente não estraguem o texto.
\IfFileExists{portuges.ldf}%
  {\usepackage[portuguese]{babel}%
   \AtBeginDocument{\shorthandoff{"}}}%
  {\usepackage[portuguese,provide=*]{babel}}
\usepackage{lmodern}       % Computer Modern em T1 (acentos nativos)
\usepackage{microtype}

\usepackage{graphicx}
\usepackage{booktabs}
\usepackage{amsmath}
\usepackage{amssymb}
\usepackage{icomma}        % vírgula decimal em modo matemático: $1,5$
\usepackage{xcolor}
\usepackage{tikz}          % diagrama da arquitetura
\usetikzlibrary{positioning,arrows.meta}

% --- Cores -------------------------------------------------------
\definecolor{corEscura}{HTML}{1F3A5F}
\definecolor{corAviso}{HTML}{B45309}

% --- Ligações ----------------------------------------------------
\usepackage[pdfusetitle]{hyperref}
\hypersetup{
  colorlinks=true,
  linkcolor=corEscura,
  citecolor=corEscura,
  urlcolor=corEscura
}
% URL e DOI das referências no tipo de letra do texto (recomendação da
% Springer), em vez de monoespaçado.
\renewcommand\UrlFont{\rmfamily}

% O cleveref não traz nomes em português; definem-se aqui.
\usepackage[capitalise,noabbrev,nameinlink]{cleveref}
\crefname{section}{Secção}{Secções}
\crefname{subsection}{Secção}{Secções}
\crefname{figure}{Figura}{Figuras}
\crefname{table}{Tabela}{Tabelas}
\crefname{equation}{Equação}{Equações}
\newcommand{\crefrangeconjunction}{ a }
\newcommand{\crefpairconjunction}{ e }
\newcommand{\crefmiddleconjunction}{, }
\newcommand{\creflastconjunction}{ e }

% --- Textos fixos da classe llncs em português -------------------
% O babel traduz «Resumo», «Figura», «Tabela» e «Referências»; as
% palavras-chave são da classe e traduzem-se à mão.
\AtBeginDocument{%
  \renewcommand{\keywordname}{\textbf{Palavras-chave:}}%
  \renewcommand{\andname}{e}%            entre dois autores
  \renewcommand{\lastandname}{\unskip\ e}% antes do último autor
}

% --- Macros próprias ---------------------------------------------
% Informação em falta, realçada a laranja. Antes de entregar não pode
% sobrar nenhuma (o build.sh não as deteta: procurar por \preencher).
\newcommand{\preencher}[1]{\textcolor{corAviso}{\textbf{[#1]}}}

% Unidades com espaço fino e sem quebra de linha: \un{2,5}{m}.
\newcommand{\un}[2]{#1\,\mbox{#2}}

% Nomes de variáveis do código (colunas das tabelas de dados).
\newcommand{\var}[1]{\texttt{\detokenize{#1}}}
